\begin{standout}[第四章]
    基本算法\\
    Basic Algorithms
\end{standout}

\begin{frame}{大纲}
    \setbeamertemplate{section in toc}[sections numbered]
    % \begin{multicols}{2}
    \tableofcontents
    % \end{multicols}
\end{frame}

\begin{frame}{学习目标}
    \begin{itemize}
        \item 理解常用排序与查找算法的基本思想及适用条件
        \item 能分析最好、最坏情况下的时间和辅助空间复杂度
        \item 能判断排序的稳定性，按数据特点选择算法
    \end{itemize}
\end{frame}

\section{排序}

\begin{frame}
    \frametitle{\insertsectionhead：引例}
    \begin{exampleblock}{整理成绩记录}
        六条记录的成绩依次为$21,25_a,22,10,25_b,18$，
        其中$a,b$表示两位不同学生，且$a$原本排在$b$之前。
        如何按成绩升序排列，同时保留同分学生原有的次序？
    \end{exampleblock}
    \begin{block}{关键需求}
        需要重新排列完整记录，而不只是数字；同键记录的先后次序也是要求。
    \end{block}
\end{frame}

\begin{frame}
    \frametitle{\insertsectionhead}
    \begin{block}{基本概念}
        \begin{itemize}
            \item 已知数据元素序列$\{a_{i}\}_{i=0}^{n}$，为每个元素$a_{i}$均关联
                  一个\textbf{关键码(key)}\footnote{简称\textbf{键}}$k_{i}$
            \item 为关键码定义一种\textbf{顺序(order)}，使得关键码之间可互相比较
                  大小
            \item 对下标序列$\{i\}_{i=0}^{n}$重\alert{排列}生成新序列
                  $\{p_{i}\}_{i=0}^{n}$
            \item 若满足$k_{p_{0}}\leq{}k_{p_{1}}\leq\cdots\leq{}k_{p_{n}}$，则
                  称$\{a_{p_{i}}\}_{i=0}^{n}$为\textbf{有序序列(sorted
                      sequence)}
            \item 称生成有序序列的过程为\textbf{排序(sorting)}
        \end{itemize}
    \end{block}
\end{frame}

\begin{frame}
    \frametitle{\insertsectionhead}
    \begin{exampleblock}{例}
        \begin{itemize}
            \item 已知整数序列：$21, 25, 22, 10, 25, 18$
            \item \alert{关键码}为元素值本身
            \item 约定\alert{顺序}为升序
            \item 排序结果为：$10, 18, 21, 22, 25, 25$
        \end{itemize}
    \end{exampleblock}
\end{frame}

\begin{fragile}
    \frametitle{\insertsectionhead}
    \begin{block}{排序的分类方法}
        \vspace*{2ex}
        \bicolumns[0.5]{
            \begin{itemize}
                \item 根据\alert{数据规模}与\alert{存储特点}分：
                      \begin{itemize}
                          \item \textbf{内部排序}：规模较小，内存足矣
                                \only<2>{$\checkmark$}
                          \item \textbf{外部排序}：规模庞大，须借助外存
                                \footnotemark
                      \end{itemize}
                \item 根据\alert{输入形式}分：
                      \begin{itemize}
                          \item \textbf{离线排序}：一次性给出全部数据
                                \only<2>{$\checkmark$}
                          \item \textbf{在线排序}：实时生成数据，分批给出
                      \end{itemize}
            \end{itemize}
        }{
            \begin{itemize}
                \item 根据\alert{体系结构}分：
                      \begin{itemize}
                          \item \textbf{串行排序}：按顺序依次处理
                                \only<2>{$\checkmark$}
                          \item \textbf{并行排序}：部分同时处理
                      \end{itemize}
                \item 根据\alert{算法策略}分：
                      \begin{itemize}
                          \item \textbf{确定排序}：不采用随机策略
                                \only<2>{$\checkmark$}
                          \item \textbf{随机排序}：采用随机策略
                      \end{itemize}
            \end{itemize}
        }
    \end{block}

    \footnotetext{如硬盘、分布式设备等}
\end{fragile}

\begin{frame}
    \frametitle{\insertsectionhead}
    \begin{block}{排序的种类}
        \begin{table}
            \centering
            \smaller
            \begin{tabular}{l|l}
                \toprule
                \textbf{算法}                   & \textbf{特点}              \\
                \midrule
                插入类：\alert{直接插入}              & \multirow{2}{*}{顺序比较、简单} \\
                选择类：\alert{直接选择}、\alert{冒泡}、堆 &                          \\
                \midrule
                \alert{快速}、归并                 & 分治策略                     \\
                \midrule
                希尔、桶、基数                       & 其他                       \\
                \bottomrule
            \end{tabular}
        \end{table}
    \end{block}
\end{frame}

\begin{fragile}
    \frametitle{\insertsectionhead}
    \begin{block}{插入排序：整理一手乱序扑克牌}
        \begin{itemize}
            \item 将序列$s$分为\alert{无序}与\alert{有序}两个子序列$s_{u}=s$与
                  $s_{o}=\emptyset$
            \item 反复从$s_{u}$中取出元素并插入$s_{o}$使之依然有序，直至
                  $s_{u}=\emptyset$，则$s_{o}$即为所求
        \end{itemize}
    \end{block}
    \bicolumns[0.5]{
        \begin{block}{算法描述}
            \begin{algorithmx}[alg:insertion]{插入排序}
                \Procedure{InsersionSort}{$s,f,t$}
                \State \Comment{对序列$s$中下标区间$[f,t)$内的元素进行插入排序}
                \State $i\gets f$
                \While{$i<t$}
                \State $s\gets$ InsertElement($s,f,i$)\Comment{插入元素子算法}
                \State $i\gets i+1$
                \EndWhile\label{insertionendwhile}
                \State \textbf{return} $s$\Comment{返回排序后的序列$s$}
                \EndProcedure
            \end{algorithmx}
        \end{block}
    }{
        \begin{exampleblock}{实例}
            \begin{itemize}
                \item 原序列：$()\alert{21}, 25, 22, 10, 25, 18$
                \item 第一趟：$(21), \alert{25}, 22, 10, 25, 18$
                \item 第二趟：$(21, 25), \alert{22}, 10, 25, 18$
                \item 第三趟：$(21, 22, 25), \alert{10}, 25, 18$
                \item 第四趟：$(10, 21, 22, 25), \alert{25}, 18$
                \item 第五趟：$(10, 21, 22, 25, 25), \alert{18}$
                \item 第六趟：$(10, 18, 21, 22, 25, 25)$
            \end{itemize}
        \end{exampleblock}
    }
\end{fragile}

\begin{frame}
    \frametitle{\insertsectionhead}
    \begin{block}{插入元素}
        \begin{itemize}
            \item 在\alert{有序序列}中\alert{查找}并\alert{插入}指定元素
        \end{itemize}
    \end{block}
    \begin{block}{可选思路}
        \begin{itemize}
            \item 顺序查找\footnote{称对应算法为\textbf{直接插入排序}}：按前驱后
                  继顺序依次遍历子序列元素，找出目标元素的合理插入位置
            \item 一些改进：二分/折半查找、Fibonacci查找等；连续存储时仍须移动元素
        \end{itemize}
    \end{block}
    \begin{block}{存储方式}
        \begin{itemize}
            \item 连续存储：须移动大量元素
            \item 链式存储：只需移动局部元素，较方便
        \end{itemize}
    \end{block}
\end{frame}

\begin{frame}
    \frametitle{\insertsectionhead}
    \begin{alertblock}{排序算法的稳定性}
        \begin{itemize}
            \item 若键相同的元素在排序前后的相对顺序保持不变，则称该排序算法为
                  \textbf{稳定的}
            \item 稳定排序算法适用于\alert{多键排序}\footnote{先按次键排序，再
                      稳定地按主键排序；主键相同的元素保留次键顺序。}
        \end{itemize}
    \end{alertblock}
\end{frame}

\begin{frame}
    \frametitle{\insertsectionhead}
    \begin{block}{直接插入排序算法性能分析}
        \begin{itemize}
            \item 时间复杂度：最好$O(n)$，平均与最坏$O(n^{2})$
                  \begin{itemize}
                      \item 平均情况依赖输入分布及具体比较方式
                  \end{itemize}
            \item 辅助空间复杂度：$O(1)$
                  \begin{itemize}
                      \item 须额外一个元素的空间缓存待交换元素
                  \end{itemize}
            \item 稳定性
                  \begin{itemize}
                      \item 在插入元素时若遇键相同的情况可小心处理顺序，故该算法
                            \alert{稳定}
                  \end{itemize}
        \end{itemize}
    \end{block}
\end{frame}

\begin{fragile}
    \frametitle{\insertsectionhead}
    \small
    \begin{block}{选择排序：择优录取}
        \begin{itemize}
            \item 将$s$分为无序子序列$s_{u}=s$与有序子序列$s_{o}=\emptyset$
            \item 反复选取$s_{u}$中的最小元素追加至$s_{o}$，直至$s_{u}=\emptyset$
        \end{itemize}
    \end{block}
    \vspace{-1.5ex}
    \bicolumns[0.5]{
        \begin{block}{算法描述}
            \begin{algorithmx}[alg:selection]{选择排序}
                \footnotesize
                \Procedure{SelectionSort}{$s,f,t$}
                \State \Comment{对序列$s$中下标区间$[f,t)$内的元素进行选择排序}
                \State $i\gets f$
                \While{$i<t$}
                \State $k\gets$ SelectElement($s,i,t$)\Comment{选择元素子算法}
                \State Swap($s_{i},s_{k}$)\Comment{交换元素}
                \State $i\gets i+1$
                \EndWhile
                \State \textbf{return} $s$\Comment{返回排序后的序列$s$}
                \EndProcedure
            \end{algorithmx}
        \end{block}
    }{
        \begin{exampleblock}{实例}
            \footnotesize
            \begin{itemize}
                \item 原序列：$()21, 25, 22, \alert{10}, 25, 18$
                \item 第一趟：$(10), 25, 22, 21, 25, \alert{18}$
                \item 第二趟：$(10, 18), 22, \alert{21}, 25, 25$
                \item 第三趟：$(10, 18, 21), \alert{22}, 25, 25$
                \item 第四趟：$(10, 18, 21, 22), \alert{25}, 25$
                \item 第五趟：$(10, 18, 21, 22, 25), \alert{25}$
                \item 第六趟：$(10, 18, 21, 22, 25, 25)$
            \end{itemize}
        \end{exampleblock}
    }
\end{fragile}

\begin{frame}
    \frametitle{\insertsectionhead}
    \begin{block}{选择元素}
        \begin{itemize}
            \item 在\alert{无序序列}中\alert{查找}键最小元素
        \end{itemize}
    \end{block}
    \begin{block}{可选思路}
        \begin{itemize}
            \item \textbf{冒泡排序}：通过逐个交换相邻元素的方式\alert{选出}键最
                  小元素
            \item \textbf{直接选择排序}：按前驱后继顺序依次遍历子序列元
                  素，\alert{找出}键最小元素
            \item \textbf{堆排序}：采用\alert{优先队列}组织无序子序列
        \end{itemize}
    \end{block}
\end{frame}

\begin{frame}
    \frametitle{\insertsectionhead}
    \begin{block}{直接选择排序的性能分析}
        \begin{itemize}
            \item 时间复杂度：最好、最坏与平均情况均为$O(n^{2})$
            \item 辅助空间复杂度：$O(1)$
            \item 稳定性
            \begin{itemize}
                \item 因涉及到跨元素交换，故该算法\alert{不稳定}
            \end{itemize}
        \end{itemize}
    \end{block}
\end{frame}

\begin{fragile}
    \frametitle{\insertsectionhead}
    \small
    \begin{block}{快速排序：分治策略的应用}
        \begin{itemize}
            \item 将序列\alert{划分}为三部分：任意元素$s_{m}$作为\alert{支点}、
                  左侧元素不大于支点，右侧元素不小于支点
            \item 对左右子序列递归执行上述\alert{划分}策略，直至子序列长度不超过$1$
        \end{itemize}
    \end{block}
    \bicolumns[0.5]{
        \begin{block}{算法描述}
            \begin{algorithmx}[alg:quick]{快速排序}
                \Procedure{QuickSort}{$s,l,h$}
                \State \Comment{对序列$s$中下标区间$[l,h]$内的元素进行快速排序}
                \If{$l<h$}
                \State $m\gets$ Partition($s,l,h$)\Comment{序列划分子算法}
                \State QuickSort($s,l,m-1$)\Comment{左侧递归}
                \State QuickSort($s,m+1,h$)\Comment{右侧递归}
                \EndIf
                \State \textbf{return} $s$\Comment{返回排序后的序列$s$}
                \EndProcedure
            \end{algorithmx}
        \end{block}
    }{
        \begin{exampleblock}{实例}
            \begin{itemize}
                \item 原序列：$21, 25, 22, 10, 25, 18$
                \item 第一趟：$18, 10, \alert{21}, 22, 25, 25$
                \item 第二趟：$10, \alert{18}, \alert{21}, \alert{22}, 25, 25$
                \item 第三趟：$\alert{10}, \alert{18}, \alert{21}, \alert{22}, \alert{25}, 25$
                \item 第四趟：$\alert{10}, \alert{18}, \alert{21}, \alert{22}, \alert{25}, \alert{25}$
            \end{itemize}
        \end{exampleblock}
    }
\end{fragile}

\begin{fragile}
    \frametitle{\insertsectionhead}
    \small
    \bicolumns[0.5]{
        \begin{block}{划分算法}
            \begin{minted}{c}
                int partition(DataType *s, int l, int h,
                              CompareType c) {
                    DataType pivot = s[l]; // 缓存支点元素
                    for (;;) {
                        for (; l < h && c(&pivot, s + h) < 0; --h)
                            ; // 左移h指针至须交换处
                        if (l >= h) break;
                        s[l++] = s[h]; // 交换元素并右移l指针
                        for (; l < h && c(s + l, &pivot) < 0; ++l)
                            ; // 右移l指针至须交换处
                        if (l >= h) break;
                        s[h--] = s[l]; // 交换元素并左移h指针
                    } // 此时l与h二指针重合，结束循环
                    s[l] = pivot; // 放回支点元素
                    return l; // 返回支点元素新下标
                }
            \end{minted}
        \end{block}
    }{
        \begin{exampleblock}{实例：一次划分}
            \footnotesize
            \begin{itemize}
                \item 原始序列：$21_{l}, 25, 22, 10, 25, 18_{h}, (\,\,\,\,\,\,)$
                \item 缓存支点：$\alert{\,\,\,\,\,\,_{l}}, 25, 22, 10, 25, \alert{18_{h}},(21_{m})$
                \item \alert{交换元素}：$\alert{18_{l}}, 25, 22, 10, 25, \alert{\,\,\,\,\,\,_{h}},(21_{m})$
                \item 移动指针：$18, \alert{25_{l}}, 22, 10, 25, \alert{\,\,\,\,\,\,_{h}},(21_{m})$
                \item \alert{交换元素}：$18, \alert{\,\,\,\,\,\,_{l}}, 22, 10, 25, \alert{25_{h}},(21_{m})$
                \item 移动指针：$18, \alert{\,\,\,\,\,\,_{l}}, 22, 10, \alert{25_{h}}, 25,(21_{m})$
                \item 移动指针：$18, \alert{\,\,\,\,\,\,_{l}}, 22, \alert{10_{h}}, 25, 25,(21_{m})$
                \item \alert{交换元素}：$18, \alert{10_{l}}, 22, \alert{\,\,\,\,\,\,_{h}}, 25, 25,(21_{m})$
                \item 移动指针：$18, 10, \alert{22_{l}}, \alert{\,\,\,\,\,\,_{h}}, 25, 25,(21_{m})$
                \item \alert{交换元素}：$18, 10, \alert{\,\,\,\,\,\,_{l}}, \alert{22_{h}}, 25, 25,(21_{m})$
                \item 移动指针：$18, 10, \alert{\,\,\,\,\,\,_{lh}}, 22, 25, 25,(21_{m})$
                \item 放回支点：$18, 10, \alert{21_{m}}, 22, 25, 25,(\,\,\,\,\,\,)$
            \end{itemize}
        \end{exampleblock}
    }
\end{fragile}

\begin{frame}
    \frametitle{\insertsectionhead}
    \begin{block}{快速排序的性能分析}
        \begin{itemize}
            \item 时间复杂度
                  \begin{itemize}
                      \item 最好情况：$O(n\log n)$
                      \item 最坏情况：$O(n^{2})$
                      \item 平均情况：$O(n\log n)$
                  \end{itemize}
            \item 递归调用的辅助空间：平均$O(\log n)$，最坏$O(n)$
            \item 稳定性
                  \begin{itemize}
                      \item 由于交换过程会越过支点元素且无法控制，故该算法
                            \alert{不稳定}
                  \end{itemize}
        \end{itemize}
    \end{block}
\end{frame}

\begin{frame}
    \frametitle{\insertsectionhead}
    \begin{alertblock}{各种排序算法的取舍}
        \begin{itemize}
            \item 问题规模较小：直接选择排序、直接插入排序等；具体阈值取决于
                  数据、实现与运行环境
            \item 问题规模略大：快速排序、归并排序、堆排序等
            \item 当需要多键排序时，宜选用稳定算法：直接插入排序、归并排序等
        \end{itemize}
    \end{alertblock}
\end{frame}

\begin{fragile}
    \frametitle{\insertsectionhead}
    \begin{block}{归并排序：分治策略的应用}
        将序列分为两半，分别递归排序，再将两个有序子序列合并。
    \end{block}
    \bicolumns[0.5]{
        \begin{block}{算法描述}
            \begin{algorithmx}[alg:merge_sort]{归并排序}
                \Procedure{MergeSort}{$s,l,r$}
                \State \Comment{排序左闭右开区间$[l,r)$}
                \If{$r-l\leq 1$}
                \State \textbf{return} $s[l:r]$
                \EndIf
                \State $m\gets l+\lfloor(r-l)/2\rfloor$
                \State $a\gets$ MergeSort($s,l,m$)
                \State $b\gets$ MergeSort($s,m,r$)
                \State \textbf{return} Merge($a,b$)
                \EndProcedure
            \end{algorithmx}
        \end{block}
    }{
        \begin{exampleblock}{实例}
            \begin{itemize}
                \item 原序列：$21,25,22,10,25,18$
                \item 分成：$(21,25,22)$、$(10,25,18)$
                \item 递归排序：$(21,22,25)$、$(10,18,25)$
                \item 合并结果：$10,18,21,22,25,25$
            \end{itemize}
        \end{exampleblock}
    }
\end{fragile}

\begin{frame}
    \frametitle{\insertsectionhead}
    \begin{block}{合并两个有序序列}
        \begin{itemize}
            \item 比较两个序列的当前元素，将较小者依次写入辅助数组
            \item 键相同时先取左侧元素，可保持相同键的原有相对顺序
            \item 一个序列耗尽后，复制另一个序列的剩余元素
        \end{itemize}
    \end{block}
    \begin{exampleblock}{实例：最后一次合并}
        $(21,22,25_a)$与$(10,18,25_b)$依次取出
        $10,18,21,22,25_a,25_b$；相等时先取左侧$25_a$。
    \end{exampleblock}
    \begin{alertblock}{性能}
        \begin{itemize}
            \item 时间：最好、平均和最坏均为$O(n\log n)$
            \item 常见数组实现的辅助空间为$O(n)$；上述合并规则保证稳定性
        \end{itemize}
    \end{alertblock}
\end{frame}

\begin{frame}
    \frametitle{\insertsectionhead}
    \begin{table}
        \centering
        \small
        \caption{常见数组排序算法比较}
        \begin{tabular}{lcccc}
            \toprule
            算法 & 最好时间 & 最坏时间 & 辅助空间 & 稳定 \\
            \midrule
            直接插入 & $O(n)$ & $O(n^2)$ & $O(1)$ & 是 \\
            直接选择 & $O(n^2)$ & $O(n^2)$ & $O(1)$ & 否 \\
            快速排序 & $O(n\log n)$ & $O(n^2)$ & $O(n)$ & 否 \\
            归并排序 & $O(n\log n)$ & $O(n\log n)$ & $O(n)$ & 是 \\
            \bottomrule
        \end{tabular}
    \end{table}
    \small 快速排序表中空间为所示递归实现的最坏情况；稳定性取决于具体实现。
\end{frame}

\begin{frame}
    \frametitle{\insertsectionhead：回到成绩记录}
    \begin{exampleblock}{选择稳定排序}
        对$21,25_a,22,10,25_b,18$采用示例中的稳定归并规则，
        得$10,18,21,22,25_a,25_b$；两位同分学生的原有次序得以保留。
    \end{exampleblock}
    \begin{block}{实际取舍}
        每次移动的是整条学生记录；归并排序提供$O(n\log n)$最坏时间，
        但常见数组实现需要$O(n)$辅助空间。
    \end{block}
\end{frame}

\begin{frame}{排序：思考与自测}
    \begin{itemize}
        \item 对$(2_a,1,2_b)$使用交换式直接选择排序，相同键的次序是否一定保持？
        \item 以成绩为主键、姓名为次键升序排列记录，如何利用稳定排序？
    \end{itemize}
    \pause
    \begin{block}<2-|handout:0>{参考答案}
        第一趟交换后为$(1,2_a,2_b)$；但$(2_a,2_b,1)$会变为
        $(1,2_b,2_a)$，故不稳定。多键排序先按姓名，再稳定地按成绩排序。
    \end{block}
\end{frame}

\section{查找}

\begin{frame}
    \frametitle{\insertsectionhead：引例}
    \begin{exampleblock}{查询图书编号}
        一批图书编号为$10,18,21,25,32$。需要判断编号$25$是否存在，
        有时还要列出编号在$18$至$25$之间的图书。
        如果查询频繁，应如何组织和查找这些记录？
    \end{exampleblock}
    \begin{block}{关键需求}
        精确查找与范围查询的要求不同，还要考虑数据是否有序以及更新代价。
    \end{block}
\end{frame}

\begin{frame}
    \frametitle{\insertsectionhead}
    \begin{block}{查找(Search)}
        \begin{itemize}
            \item 根据给定键定位记录，或判断记录不存在
            \item 顺序查找依次比较元素；无需有序，最坏时间为$O(n)$
            \item 查找算法的选择取决于数据是否有序、存储方式及更新频率
        \end{itemize}
    \end{block}
\end{frame}

\begin{frame}
    \frametitle{\insertsectionhead}
    \begin{block}{二分查找(Binary Search)}
        \begin{itemize}
            \item 前提：序列按键有序，且可高效随机访问
            \item 每次比较中间元素，将候选区间缩小至一半
            \item 以下使用左闭右开区间$[l,r)$，未找到时返回$-1$
        \end{itemize}
    \end{block}
\end{frame}

\begin{fragile}
    \frametitle{\insertsectionhead}
    \begin{algorithmx}[alg:binary-search]{二分查找}
        \Procedure{BinarySearch}{$s,n,key$}
        \State $l\gets 0,\ r\gets n$
        \While{$l<r$}
        \State $m\gets l+\lfloor(r-l)/2\rfloor$
        \If{$s_m=key$}
        \State \textbf{return} $m$
        \ElsIf{$s_m<key$}
        \State $l\gets m+1$
        \Else
        \State $r\gets m$
        \EndIf
        \EndWhile
        \State \textbf{return} $-1$
        \EndProcedure
    \end{algorithmx}
\end{fragile}

\begin{frame}
    \frametitle{\insertsectionhead}
    \begin{exampleblock}{二分查找实例}
        在$(10,18,21,25,32)$中查找$25$：
        \[
            [0,5)\xrightarrow{21<25}[3,5)
            \xrightarrow{32>25}[3,4)
            \xrightarrow{25=25}\text{找到下标}3
        \]
    \end{exampleblock}
    \begin{itemize}
        \item 最好时间$O(1)$，最坏时间$O(\log n)$；迭代实现辅助空间$O(1)$
        \item 键重复时，该算法返回任意匹配位置；首末位置需调整边界规则
        \item 空序列与查找失败时，区间为空，返回$-1$
    \end{itemize}
\end{frame}

\begin{frame}
    \frametitle{\insertsectionhead}
    \begin{block}{散列表(Hash Table)}
        \begin{itemize}
            \item 散列函数将键映射到表中位置；不同键可能映射到同一位置，称为冲突
            \item 冲突可用链地址法或开放定址法处理；查找和插入需使用一致的规则
            \item 在散列分布良好且装载因子受控时，查找平均为$O(1)$；最坏可为$O(n)$
            \item 普通散列表不保持键的顺序，不适合直接进行范围查询
        \end{itemize}
    \end{block}
    \begin{exampleblock}{链地址法示例}
        表长$5$，$h(k)=k\bmod 5$：键$12$与$17$均落在桶$2$，
        在该桶中依次比较即可区分。
    \end{exampleblock}
\end{frame}

\begin{frame}
    \frametitle{\insertsectionhead}
    \begin{block}{散列函数与装载因子}
        \begin{itemize}
            \item 散列函数$h(k)$把键$k$映射到表下标，计算应快速且分布尽量均匀
            \item 除留余数法示例：表长$m=7$，$h(k)=k\bmod 7$
            \item 已存储$n$个键时，装载因子$\alpha=n/m$反映表的拥挤程度
            \item 装载因子增加通常会增加冲突；可扩容并重新散列所有键
        \end{itemize}
    \end{block}
\end{frame}

\begin{frame}
    \frametitle{\insertsectionhead}
    \begin{block}{链地址法(Separate Chaining)}
        \begin{itemize}
            \item 每个桶保存映射到该位置的键；冲突键放入同一链表
            \item 插入、查找和删除先定位桶，再在该桶中查找目标键
            \item 均匀散列假设下平均查找时间为$O(1+\alpha)$；最坏为$O(n)$
            \item 桶中可存多个键，因此$\alpha$可以大于$1$
        \end{itemize}
    \end{block}
    \begin{exampleblock}{例：$m=7$，$h(k)=k\bmod 7$}
        插入$10,17,24$后，桶$3$的链表可为$10\rightarrow17\rightarrow24$。
        查找$17$需在桶$3$中比较，查找$31$遍历该桶后失败。
    \end{exampleblock}
\end{frame}

\begin{frame}
    \frametitle{\insertsectionhead}
    \begin{block}{开放定址法(Open Addressing)：线性探测}
        \begin{itemize}
            \item 全部键存于表内；冲突时依次探测$(h(k)+i)\bmod m$
            \item 插入时选第一个可用位置；查找沿相同探测序列进行
            \item 找到目标或遇到从未使用过的空槽时结束；最多探测$m$次
        \end{itemize}
    \end{block}
    \begin{exampleblock}{例：$m=7$，$h(k)=k\bmod 7$}
        依次插入$10,17,24$，分别存入下标$3,4,5$；查找$24$依次检查$3,4,5$。
    \end{exampleblock}
\end{frame}

\begin{frame}
    \frametitle{\insertsectionhead}
    \begin{block}{删除与性能}
        \begin{itemize}
            \item 开放定址删除时须留下\textbf{删除标记}，不能直接改为从未使用的空槽，
                  否则可能截断后续键的查找路径
            \item 插入可复用删除槽，但若需避免重复键，应继续探测至找到该键或确认不存在
            \item 开放定址的$\alpha<1$；表接近满时探测会变长，需及时扩容并重建
            \item 在良好的散列与受控装载因子下，平均查找可为$O(1)$；最坏为$O(n)$
        \end{itemize}
    \end{block}
    \begin{alertblock}{思考}
        上例删除$17$后，为什么不能把下标$4$直接标为从未使用？
    \end{alertblock}
    \begin{block}<2-|handout:0>{参考答案}
        查找$24$会在下标$4$提前结束，无法到达下标$5$。
    \end{block}
\end{frame}


\begin{frame}
    \frametitle{\insertsectionhead：回到图书编号查询}
    \begin{exampleblock}{选择查找方式}
        \begin{itemize}
            \item 有序数组中二分查找$25$：先比较$21$，再比较$32$，最后找到$25$
            \item 表长为$5$、$h(k)=k\bmod5$的链地址散列表中，
                  $10$和$25$同在桶$0$；查找$25$须在该桶中比较
            \item 范围$[18,25]$：在有序数组中找到起点并顺序扫描，
                  输出$18,21,25$；普通散列表不保持顺序
        \end{itemize}
    \end{exampleblock}
    有序数组适合随机访问，但插入删除可能移动元素；散列表适合精确查找。
\end{frame}

\begin{frame}{查找：思考与自测}
    \begin{itemize}
        \item 在$(2,4,6,8)$中二分查找$5$，候选区间如何变化？
        \item 数据经常更新且需范围查询时，散列表是否总是合适？
    \end{itemize}
    \pause
    \begin{block}<2-|handout:0>{参考答案}
        $[0,4)\rightarrow[0,2)\rightarrow[2,2)$，未找到。
        普通散列表不维护顺序；范围查询可考虑有序结构并权衡更新成本。
    \end{block}
\end{frame}


\section{小结}

\begin{frame}
    \frametitle{\insertsectionhead}
    \begin{itemize}
        \item 排序算法需综合考虑时间、辅助空间与稳定性
        \item 二分查找依赖有序序列；散列表通过散列函数和冲突处理支持查找
        \item 复杂度分析应说明输入条件及最好、最坏或平均情况
    \end{itemize}
\end{frame}

\begin{standout}[]
    问与答
\end{standout}